% Extending the Bridge: Conditioning and Update between Probability
% and Many-Valued Logic
%
% Part 2 of the product De Morgan triplet series.
% Lean formalization: lean/Cred/{Basic,Congruence,Consequence,
%   Valuation,Update,Predicate,CondBridge}.lean

\documentclass[11pt,a4paper]{article}

\usepackage{amsmath,amssymb,amsthm}
\usepackage{mathtools}
\usepackage{stmaryrd}
\usepackage{enumitem}
\usepackage{booktabs}
\usepackage{caption}
\usepackage{hyperref}
\hypersetup{
    colorlinks=true,
    linkcolor=blue,
    citecolor=blue,
    urlcolor=blue,
    pdfauthor={Anonymous},
    pdftitle={Extending the Bridge: Conditioning and Update between Probability and Many-Valued Logic}
}
\usepackage{cleveref}
\usepackage[margin=1in]{geometry}
\usepackage{tikz}
\usetikzlibrary{arrows.meta,positioning,shapes}
\usepackage{graphicx}

% Theorem environments
\newtheorem{theorem}{Theorem}[section]
\newtheorem{lemma}[theorem]{Lemma}
\newtheorem{proposition}[theorem]{Proposition}
\newtheorem{corollary}[theorem]{Corollary}
\newtheorem{definition}[theorem]{Definition}
\newtheorem{example}[theorem]{Example}
\newtheorem{remark}[theorem]{Remark}

% Custom commands
\newcommand{\cred}{\mathsf{cred}}
\newcommand{\negc}{\mathord{\sim}}
\newcommand{\conjc}{\otimes}
\newcommand{\disjc}{\sqcup}
\newcommand{\condc}[2]{({#1} \mid {#2})}
\newcommand{\Real}{\mathbb{R}}
\newcommand{\Nat}{\mathbb{N}}

\title{Extending the Bridge: Conditioning and Update\\
between Probability and Many-Valued Logic}

\author{Anonymous for blind review}
\date{February 2026}

\begin{document}

\maketitle

\begin{abstract}
Part~1 of this series established a bridge between continuous
$[0,1]$ reasoning and three-valued logic:
the product De Morgan triplet collapses to the Kleene lattice
$\{0, \tfrac{1}{2}, 1\}$, and LP/K$_3$ consequence is equivalent
to positivity/certainty consequence on $[0,1]$, at the level of
propositional formulas.
This paper asks whether the bridge extends to conditioning.
We prove that it does not: no truth-functional function
$f : \{0, \tfrac{1}{2}, 1\}^2 \to \{0, \tfrac{1}{2}, 1\}$
can satisfy $\mathrm{collapse}(\min(a,b)/b) = f(\mathrm{collapse}(b),
\mathrm{collapse}(a))$ for all $a, b$ with $b > 0$
--- a constructive Lewis triviality result for the product
De Morgan triplet.
We characterize the boundary conditions under which the bridge
does hold (evidence certain, conclusion impossible or certain, or
conclusion at least as credible as evidence), and show that
Bayesian update inherits the bridge under these conditions.
Zero evidence ties the paper together: the algebraic fact
$c \cdot 0 = 0$ simultaneously forces underdetermined posteriors,
paraconsistent explosion failure, and conditional bridge failure.
All results are machine-checked in Lean~4 with Mathlib
(3581 lines across seven modules, zero \texttt{sorry}).
\end{abstract}

%=============================================================================
\section{Introduction}
\label{sec:intro}
%=============================================================================

You believe it will rain with credence $0.8$.
You see the street is wet with credence $0.9$.
What is your conditional credence in rain given wet streets?

In Part~1 we studied the product De Morgan triplet on $[0,1]$ ---
product conjunction ($c_1 \conjc c_2 = c_1 \cdot c_2$),
probabilistic-sum disjunction
($c_1 \disjc c_2 = c_1 + c_2 - c_1 c_2$), and standard complement
($\negc c = 1 - c$) --- and showed that the triplet collapses to
the Kleene lattice $\{0, \tfrac{1}{2}, 1\}$ via a surjective
homomorphism.
The main result of Part~1 was the \emph{bridge theorem}: LP
consequence on $\{0, \tfrac{1}{2}, 1\}$ is equivalent to positivity
consequence on $[0,1]$, and K$_3$ consequence is equivalent to
certainty consequence, both at the level of propositional formulas
built from negation, conjunction, and disjunction.

Part~1 noted, but did not pursue, a fundamental limitation: the
bridge covers connectives $(\negc, \conjc, \disjc)$, not
conditioning.
This paper takes up that question directly.

\textbf{The question.}
The collapse homomorphism preserves negation, conjunction, and
disjunction.
Does it also preserve conditioning?
More precisely: if we compute the min-copula conditional
$\min(a,b)/b$ on $[0,1]$ and then collapse, do we get the same
answer as first collapsing and then computing the G\"odel
conditional on $\{0, \tfrac{1}{2}, 1\}$?

In our running example:
$\cred(\text{rain}) = 0.8$, $\cred(\text{wet}) = 0.9$,
so $\cred(\text{rain} \mid \text{wet}) = \min(0.8, 0.9)/0.9
= 0.8/0.9 = 8/9$.
Since $8/9 \in (0,1)$, it collapses to $\tfrac{1}{2}$.
But the G\"odel conditional on the collapsed values is
$\text{g\"odel}(\tfrac{1}{2}, \tfrac{1}{2}) = 1 \neq \tfrac{1}{2}$.
The bridge is already broken.

\textbf{Main results.}
\begin{enumerate}[nosep]
\item \textbf{Propositional bridge} (Section~\ref{sec:bridge}):
  We recall the bridge theorem from Part~1 and state the
  formula-level results we need: collapse commutes with evaluation,
  LP $\leftrightarrow$ positivity, K$_3$ $\leftrightarrow$ certainty.
  These cover the connective language.
\item \textbf{Conditional bridge impossibility}
  (Section~\ref{sec:cond-bridge}):
  No truth-functional conditional bridge exists
  (Theorem~\ref{thm:no-cond-bridge}); the bridge holds exactly at
  boundary cases
  (Theorem~\ref{thm:boundary}); the bridge fails for
  all interior pairs with $a < b$
  (Theorem~\ref{thm:interior-fail}).
  This is a constructive Lewis triviality result: the impossibility
  is witnessed by two explicit pairs, not by a measure-theoretic
  argument.
\item \textbf{Update bridge}
  (Section~\ref{sec:update}):
  Bayesian update inherits the conditional bridge under the boundary
  conditions of Theorem~\ref{thm:boundary}
  (Theorem~\ref{thm:update-bridge}); qualitative update on
  $\{0, \tfrac{1}{2}, 1\}$ via G\"odel implication
  provides the discrete counterpart.
\item \textbf{Zero-evidence duality}
  (Section~\ref{sec:update}):
  The algebraic fact $c \cdot 0 = 0$ has three equivalent
  manifestations: underdetermined posteriors, LP explosion failure,
  and conditional bridge failure
  (Theorem~\ref{thm:zero-duality}).
\end{enumerate}

\textbf{Literature positioning.}
Lewis~\cite{lewis1976} showed that identifying $P(A \mid B)$ with
the probability of a conditional proposition forces triviality;
Halpern~\cite{halpern2003} analyzed the scope of Lewis's conditions.
Our impossibility theorem (Theorem~\ref{thm:no-cond-bridge}) is a
different kind of triviality: it shows the collapse homomorphism
cannot extend to conditioning, even without treating conditionals
as propositions.
Adams~\cite{adams1975} developed a probabilistic logic of
conditionals; H\'ajek~\cite{hajek1998} and
Esteva--Godo~\cite{esteva2000,esteva2001} developed Product Logic
with residuated implication (forcing ex falso);
\'Egr\'e, Rossi, and Sprenger~\cite{egre2021} developed
de Finetti's conditional algebraically.
Chain-rule conditioning~\cite{renyi1955,popper1959} avoids both
division (probability) and residuation (logic), leaving the
conditional unconstrained at zero evidence.

\textbf{Note on formalization.}
All results marked with ``Lean:'' are machine-checked in Lean~4 with
Mathlib (Appendix~\ref{app:lean}).

\textbf{Section guide.}
Section~\ref{sec:setup}: setup and notation from Part~1.
Section~\ref{sec:bridge}: the propositional bridge (warmup).
Section~\ref{sec:cond-bridge}: the conditional bridge (centerpiece).
Section~\ref{sec:update}: update through the bridge.
Section~\ref{sec:conclusion}: conclusion and open problems.

%=============================================================================
\section{Setup}
\label{sec:setup}
%=============================================================================

We work with the product De Morgan triplet on $[0,1]$ as defined in
Part~1.
This section recalls the ingredients we need; all proofs are in
Part~1 or in the Lean formalization.

\subsection{The Product De Morgan Triplet}

A \emph{credence} is a value $c \in [0,1]$.
The operations are:
negation $\negc c = 1 - c$,
conjunction $c_1 \conjc c_2 = c_1 \cdot c_2$ (the product t-norm),
and disjunction $c_1 \disjc c_2 = c_1 + c_2 - c_1 c_2$ (the
De Morgan dual).
The algebraic properties --- involution, commutativity, associativity,
De Morgan laws --- hold by arithmetic on $[0,1]$ (Part~1, Section~2).

\begin{example}[Running Example: Rain and Wet Streets]
\label{ex:rain}
Throughout this paper:
$\cred(\text{rain}) = 0.8$ and $\cred(\text{wet}) = 0.9$.
Then $\cred(\text{no rain}) = 0.2$ and
$\cred(\text{rain}) \conjc \cred(\text{wet}) = 0.72$.
\end{example}

\subsection{Chain-Rule Conditioning}

The conditional satisfies the chain rule:
$\condc{A}{B} \conjc \cred(B) = \cred(A \land B)$.
When $\cred(B) > 0$, the conditional is unique:
$\condc{A}{B} = \cred(A \land B) / \cred(B)$.
When $\cred(B) = 0$, the equation $c \cdot 0 = 0$ holds for all
$c \in [0,1]$: the conditional is unconstrained
(Lean: \texttt{conditioning\_zero\_any},
\texttt{conditioning\_unique}).
The joint $\cred(A \land B)$ is supplied as input, not computed from
marginals; the conditional depends on the dependence structure.

\subsection{Min-Copula Uniqueness}

Among symmetric copulas on $[0,1]$, the minimum is the unique
idempotent one: $C(a,b) = \min(a,b)$ is the only copula satisfying
$C(a,a) = a$ for all $a$
(Part~1, Theorem~2.7;
Lean: \texttt{copula\_idempotent\_unique},
\texttt{min\_copula\_unique}).
When the joint is the min copula, the conditional for $b > 0$ is
$\min(a,b)/b$: this equals $a/b$ when $a \leq b$, and $1$ when
$a \geq b$.

\subsection{The Collapse Homomorphism}

Define:
\[
\mathrm{collapse}(c) = \begin{cases}
0 & \text{if } c = 0, \\
\tfrac{1}{2} & \text{if } 0 < c < 1, \\
1 & \text{if } c = 1.
\end{cases}
\]
The collapse is a surjective homomorphism from
$([0,1], \negc, \conjc, \disjc)$ to
$(\{0, \tfrac{1}{2}, 1\}, \mathrm{neg}, \min, \max)$
(Part~1;
Lean: \texttt{collapse\_neg}, \texttt{collapse\_conj},
\texttt{collapse\_disj}, \texttt{collapse\_surjective}).
The target is the Kleene lattice~\cite{kleene1952}, underlying
K$_3$~\cite{kleene1952}, LP~\cite{asenjo1966,priest1979}, and
RM3~\cite{anderson1975,meyer1972}.

\subsection{The G\"odel Conditional on $\{0, \tfrac{1}{2}, 1\}$}

The product-residuated implication on $\{0, \tfrac{1}{2}, 1\}$ is
$a \to b = \sup\{c : a \cdot c \leq b\}$, which coincides with
the G\"odel implication on the three-element set
(Lean: \texttt{godel\_impl\_eq\_prod\_resid}).
Its table:
\begin{center}
\small
\begin{tabular}{r|ccc}
$a \!\downarrow\;\; b\to$ & 0 & $\tfrac{1}{2}$ & 1 \\
\hline
0 & 1 & 1 & 1 \\
$\tfrac{1}{2}$ & 0 & 1 & 1 \\
1 & 0 & $\tfrac{1}{2}$ & 1
\end{tabular}
\normalsize
\end{center}
\smallskip
\noindent
Lean: \texttt{godel\_impl},
\texttt{godel\_impl\_zero\_zero} through
\texttt{godel\_impl\_one\_one}.

\subsection{Propositional Formulas}

Propositional formulas over atoms $\alpha$ are built from variables,
negation, conjunction, and disjunction.
Given an assignment $v : \alpha \to [0,1]$, evaluation is recursive:
$\llbracket \negc\varphi \rrbracket_v = \negc\llbracket\varphi\rrbracket_v$,
$\llbracket \varphi \conjc \psi \rrbracket_v =
\llbracket\varphi\rrbracket_v \conjc \llbracket\psi\rrbracket_v$,
$\llbracket \varphi \disjc \psi \rrbracket_v =
\llbracket\varphi\rrbracket_v \disjc \llbracket\psi\rrbracket_v$.
Analogously for three-valued evaluation using Kleene connectives.
Lean: \texttt{Formula}, \texttt{evalCred}, \texttt{evalThreeVal}.

%=============================================================================
\section{The Propositional Bridge}
\label{sec:bridge}
%=============================================================================

This section summarizes the bridge theorem from Part~1 at the level
we need for the conditional bridge analysis.
All proofs are in Part~1; the Lean formalizations are in
\texttt{Cred/Consequence.lean}.

\subsection{Collapse Commutes with Evaluation}

\begin{proposition}[Collapse-Evaluation Commutativity; Part~1]
\label{prop:collapse-eval}
For any $v : \alpha \to [0,1]$ and formula $\varphi$:
\[
\mathrm{collapse}(\llbracket\varphi\rrbracket_v)
  = \llbracket\varphi\rrbracket_{\mathrm{collapse} \circ\, v}.
\]
Conversely, for any $w : \alpha \to \{0, \tfrac{1}{2}, 1\}$:
\[
\llbracket\varphi\rrbracket_w
  = \mathrm{collapse}(\llbracket\varphi\rrbracket_{\mathrm{lift} \circ\, w}).
\]
Lean: \texttt{collapse\_eval\_eq}, \texttt{lift\_eval\_eq}.
\end{proposition}

The commutativity reduces compound formulas to the atom case: the
collapse is a homomorphism of connectives, so it commutes with any
formula built from those connectives.

\subsection{Bridge Theorems}

\begin{theorem}[Formula-Level LP Bridge; Part~1]
\label{thm:lp-formula-bridge}
LP formula consequence
$\Gamma \models_{\mathrm{LP}} \varphi$
is equivalent to formula positivity consequence
$\Gamma \models^+ \varphi$:
\[
\Gamma \models_{\mathrm{LP}} \varphi
\quad \Longleftrightarrow \quad
\Gamma \models^+ \varphi.
\]
Lean: \texttt{lp\_formula\_bridge}.
\end{theorem}

\begin{theorem}[Formula-Level K$_3$ Bridge; Part~1]
\label{thm:k3-formula-bridge}
K$_3$ formula consequence
$\Gamma \models_{\mathrm{K}_3} \varphi$
is equivalent to formula certainty consequence
$\Gamma \models^1 \varphi$:
\[
\Gamma \models_{\mathrm{K}_3} \varphi
\quad \Longleftrightarrow \quad
\Gamma \models^1 \varphi.
\]
Lean: \texttt{k3\_formula\_bridge}.
\end{theorem}

\begin{corollary}[Finite Model Property; Part~1]
\label{cor:finite-model}
Positivity and certainty consequence on $[0,1]$ are decidable:
non-consequence reduces to a counterexample in
$\{0, \tfrac{1}{2}, 1\}^{|\alpha|}$.
LP and K$_3$ consequence are
coNP-complete~\cite{priest2008,avron1991}; the bridge transfers
this to positivity and certainty consequence on $[0,1]$.
\end{corollary}

\subsection{Structural Rules}

All consequence relations satisfy the standard Tarskian structural
rules --- reflexivity, monotonicity, and cut --- because they
quantify over all assignments; adding premises restricts the set of
admissible valuations
(Lean: \texttt{formulaLP\_reflexivity/monotonicity/cut},
\texttt{formulaK3\_reflexivity/monotonicity/cut},
\texttt{formulaPositivity\_reflexivity/monotonicity/cut},
\texttt{formulaCertainty\_reflexivity/monotonicity/cut}).

\subsection{Explosion Failure}
\label{sec:explosion}

\begin{example}[Explosion Failure via the Bridge]
\label{ex:explosion}
Set $\cred(\text{rain}) = \tfrac{1}{2}$ (maximum uncertainty)
and $\cred(\text{snow}) = 0$.
The positivity witness:
$v(\text{rain}) = \tfrac{1}{2} > 0$ and
$\negc v(\text{rain}) = \tfrac{1}{2} > 0$, but
$v(\text{snow}) = 0$.
Contradictory evidence about rain does not force belief in snow.
Via the bridge, the LP counterexample
$w(\text{rain}) = \tfrac{1}{2}$, $w(\text{snow}) = 0$ refutes
$\text{rain}, \negc\text{rain} \models_{\mathrm{LP}} \text{snow}$
(Lean: \texttt{lp\_no\_explosion}).
\end{example}

The bridge covers negation, conjunction, and disjunction.
Does it extend to conditioning?

%=============================================================================
\section{The Conditional Bridge}
\label{sec:cond-bridge}
%=============================================================================

The propositional bridge transfers consequence between $[0,1]$ and
$\{0, \tfrac{1}{2}, 1\}$ for formulas built from connectives.
This section asks whether the same transfer works for conditioning:
does $\mathrm{collapse}(\min(a,b)/b)$ equal a function of
$\mathrm{collapse}(a)$ and $\mathrm{collapse}(b)$?

\subsection{Min-Copula Conditioning on $[0,1]$ vs.\
$\{0, \tfrac{1}{2}, 1\}$}

Under the min copula, the conditional for $b > 0$ is
$\min(a,b)/b$.
On the Kleene lattice, the G\"odel implication (which coincides
with the product-residuated implication;
Lean: \texttt{godel\_impl\_eq\_prod\_resid}) is the natural
candidate for a three-valued conditional.

Three conditioning tables side by side, on
$\{0, \tfrac{1}{2}, 1\}$:

\begin{center}
\small
\begin{tabular}{r|ccc}
\multicolumn{4}{c}{\textbf{Kleene material}\textsuperscript{1}: $\max(\negc e, p)$} \\
$e \!\downarrow\;\; p$ & 0 & $\tfrac{1}{2}$ & 1 \\
\hline
0 & 1 & 1 & 1 \\
$\tfrac{1}{2}$ & $\tfrac{1}{2}$ & $\tfrac{1}{2}$ & 1 \\
1 & 0 & $\tfrac{1}{2}$ & 1
\end{tabular}
\quad
\begin{tabular}{r|ccc}
\multicolumn{4}{c}{\textbf{Prod.\ resid.}\textsuperscript{2}: $\min(p/e, 1)$} \\
$e \!\downarrow\;\; p$ & 0 & $\tfrac{1}{2}$ & 1 \\
\hline
0 & 1 & 1 & 1 \\
$\tfrac{1}{2}$ & 0 & 1 & 1 \\
1 & 0 & $\tfrac{1}{2}$ & 1
\end{tabular}
\quad
\begin{tabular}{r|ccc}
\multicolumn{4}{c}{\textbf{Chain-rule with min-copula joint}} \\
$e \!\downarrow\;\; p$ & 0 & $\tfrac{1}{2}$ & 1 \\
\hline
0 & $*$ & $*$ & $*$ \\
$\tfrac{1}{2}$ & 0 & 1 & 1 \\
1 & 0 & $\tfrac{1}{2}$ & 1
\end{tabular}
\normalsize
\end{center}
\smallskip
\noindent
{\footnotesize Rows = evidence $e$, columns = conclusion $p$.
$*$ = unconstrained (any value in $[0,1]$ satisfies the chain rule).
\textsuperscript{1}Shared by K$_3$~\cite{kleene1952},
LP~\cite{asenjo1966,priest1979}, RM3~\cite{anderson1975}.
\textsuperscript{2}Coincides with G\"odel on $\{0, \tfrac{1}{2}, 1\}$
(Lean: \texttt{godel\_impl\_eq\_prod\_resid}).}

\begin{example}[Running Example: Bridge Failure]
\label{ex:bridge-failure}
$\cred(\text{rain}) = 0.8$, $\cred(\text{wet}) = 0.9$.
The min-copula conditional:
$\cred(\text{rain} \mid \text{wet})
= \min(0.8, 0.9)/0.9 = 8/9 \approx 0.889$.
Collapsing: $\mathrm{collapse}(8/9) = \tfrac{1}{2}$ (interior value).
G\"odel on collapsed values:
$\mathrm{g\ddot{o}del}(\mathrm{collapse}(0.9),
\mathrm{collapse}(0.8))
= \mathrm{g\ddot{o}del}(\tfrac{1}{2}, \tfrac{1}{2}) = 1$.
The two paths disagree: $\tfrac{1}{2} \neq 1$.
\end{example}

\subsection{Impossibility Theorem}

\begin{theorem}[No Truth-Functional Conditional Bridge; new]
\label{thm:no-cond-bridge}
There is no function
$f : \{0, \tfrac{1}{2}, 1\}^2 \to \{0, \tfrac{1}{2}, 1\}$
satisfying
\[
\mathrm{collapse}\!\left(\frac{\min(a,b)}{b}\right)
= f(\mathrm{collapse}(b),\, \mathrm{collapse}(a))
\]
for all $a, b \in [0,1]$ with $b > 0$.

\smallskip\noindent
Lean: \texttt{no\_truthfunctional\_cond\_bridge}.
\end{theorem}

\begin{proof}
Two witnesses force contradictory values for
$f(\tfrac{1}{2}, \tfrac{1}{2})$.

\emph{Witness 1:}
$a = 0.3$, $b = 0.7$.
Both are interior, so
$\mathrm{collapse}(a) = \mathrm{collapse}(b) = \tfrac{1}{2}$.
The conditional is $\min(0.3, 0.7)/0.7 = 0.3/0.7 = 3/7 \in (0,1)$,
so $\mathrm{collapse}(3/7) = \tfrac{1}{2}$.
This forces $f(\tfrac{1}{2}, \tfrac{1}{2}) = \tfrac{1}{2}$.

\emph{Witness 2:}
$a = 0.7$, $b = 0.3$.
Again $\mathrm{collapse}(a) = \mathrm{collapse}(b) = \tfrac{1}{2}$.
The conditional is $\min(0.7, 0.3)/0.3 = 0.3/0.3 = 1$,
so $\mathrm{collapse}(1) = 1$.
This forces $f(\tfrac{1}{2}, \tfrac{1}{2}) = 1$.

Contradiction: $\tfrac{1}{2} \neq 1$.
\end{proof}

\begin{remark}[Constructive Lewis triviality]
\label{rem:lewis}
Lewis~\cite{lewis1976} showed that identifying $P(A \mid B)$ with
the probability of a conditional proposition $A \to B$ in the same
algebra forces the probability measure to be trivial.
Theorem~\ref{thm:no-cond-bridge} is a different kind of triviality:
it shows the collapse homomorphism, which faithfully preserves
connectives, \emph{cannot} extend to conditioning.
The obstacle is not that conditionals are propositions (they are not,
in our framework), but that the collapsed values lose the information
needed to determine the conditional.
Two pairs $(a, b)$ with the same collapsed image can have different
conditional values because $\min(a,b)/b$ depends on the ratio, not
the collapsed class.
Lewis requires measure-theoretic machinery; our result is finitary
(two explicit witnesses).
See Halpern~\cite{halpern2003} for a detailed analysis of Lewis's
conditions;
Adams~\cite{adams1975} for a probabilistic logic of conditionals
that avoids Lewis's assumptions differently.
\end{remark}

\subsection{Characterization: When the Bridge Works}

Although the conditional bridge fails in general, it holds at the
boundaries.

\begin{theorem}[Boundary Characterization; new]
\label{thm:boundary}
For $a, b \in [0,1]$ with $b > 0$:
\[
\mathrm{collapse}\!\left(\frac{\min(a,b)}{b}\right)
= \mathrm{g\ddot{o}del}(\mathrm{collapse}(b),\,
\mathrm{collapse}(a))
\]
whenever at least one of the following holds:
\begin{enumerate}[label=\textup{(\roman*)},nosep]
\item $a = 0$ (conclusion impossible),
\item $a = 1$ (conclusion certain),
\item $b = 1$ (evidence certain),
\item $a \geq b$ (conclusion at least as credible as evidence).
\end{enumerate}
In all four cases, the min-copula conditional is either $0$ or $1$
(a boundary value).

\smallskip\noindent
Lean: \texttt{cond\_bridge\_boundary},
\texttt{cond\_bridge\_at\_a\_zero},
\texttt{cond\_bridge\_at\_a\_one},
\texttt{cond\_bridge\_at\_b\_one},
\texttt{cond\_bridge\_at\_a\_ge\_b}.
\end{theorem}

\begin{proof}
Case analysis.
(i)~$a = 0$: $\min(0,b)/b = 0$; $\mathrm{collapse}(0) = 0$;
$\mathrm{g\ddot{o}del}(\mathrm{collapse}(b), 0) = 0$ for
$b > 0$.
(ii)~$a = 1$: $\min(1,b)/b = b/b = 1$;
$\mathrm{g\ddot{o}del}(\mathrm{collapse}(b), 1) = 1$ for all
$b$.
(iii)~$b = 1$: $\min(a,1)/1 = a$; collapsing commutes trivially.
(iv)~$a \geq b$: $\min(a,b)/b = b/b = 1$;
$\mathrm{collapse}(1) = 1$; and
$\mathrm{g\ddot{o}del}(\mathrm{collapse}(b),
\mathrm{collapse}(a)) = 1$ when $a \geq b > 0$.
\end{proof}

\begin{example}[Running Example: When the Bridge Works]
\label{ex:bridge-works}
If $\cred(\text{rain}) = 0.95 > 0.9 = \cred(\text{wet})$, then
$\min(0.95, 0.9)/0.9 = 1$ (case (iv): $a \geq b$).
$\mathrm{collapse}(1) = 1 =
\mathrm{g\ddot{o}del}(\tfrac{1}{2}, \tfrac{1}{2})$?
No: $\mathrm{g\ddot{o}del}(\tfrac{1}{2}, \tfrac{1}{2}) = 1$,
so in this case the bridge does hold.
The boundary condition is that one path or the other forces a
boundary conditional.
\end{example}

\subsection{Interior Failure}

The bridge fails precisely when both $a$ and $b$ are interior and
$a < b$.

\begin{theorem}[Interior Bridge Failure; new]
\label{thm:interior-fail}
There exist $a, b \in (0,1)$ with $a < b$ such that
\[
\mathrm{collapse}\!\left(\frac{\min(a,b)}{b}\right)
\neq \mathrm{g\ddot{o}del}(\mathrm{collapse}(b),\,
\mathrm{collapse}(a)).
\]
More precisely: for any interior $a, b$ with $0 < a < b < 1$, the
conditional $a/b$ is interior (strictly between $0$ and $1$), so
$\mathrm{collapse}(a/b) = \tfrac{1}{2}$; but
$\mathrm{g\ddot{o}del}(\tfrac{1}{2}, \tfrac{1}{2}) = 1$.

\smallskip\noindent
Lean: \texttt{cond\_bridge\_fails\_interior},
\texttt{minCopulaCond\_interior}.
\end{theorem}

\begin{proof}
When $0 < a < b < 1$: $\min(a,b)/b = a/b$.
Since $0 < a < b$, we have $0 < a/b < 1$, so
$\mathrm{collapse}(a/b) = \tfrac{1}{2}$.
Since $0 < a < 1$ and $0 < b < 1$, both collapse to
$\tfrac{1}{2}$, and
$\mathrm{g\ddot{o}del}(\tfrac{1}{2}, \tfrac{1}{2}) = 1
\neq \tfrac{1}{2}$.
The witness $a = 3/10$, $b = 7/10$ makes this concrete.
\end{proof}

The failure is systematic, not accidental.
Whenever both the conclusion and evidence are uncertain (interior)
and the conclusion is less credible than the evidence, the
min-copula conditional is interior; but the G\"odel conditional on
$\{0, \tfrac{1}{2}, 1\}$ maps $(\tfrac{1}{2}, \tfrac{1}{2})$ to $1$
(certainty), losing the distinction between ``partially supports''
and ``fully supports.''

\subsection{Path Dependence}
\label{sec:path-dep}

The impossibility theorem has a topological manifestation.
Under the min copula, the conditional $\min(a,b)/b$ as $b \to 0$ is
path-dependent: the limit depends on \emph{how} $a$ and $b$ approach
zero together.

\begin{figure}[ht]
\centering
\includegraphics[width=\textwidth]{figures/path_dependence}
\caption{Path dependence at evidence zero.
\textbf{Left:} $\min(a,b)/b$ near the origin; value depends on angle
$\theta$, not radius $r$.
\textbf{Right:} Limit as $r \to 0$: $1$ for $\theta \leq \pi/4$
(since $a \geq b$ gives $\min(a,b)/b = 1$),
$\cot\theta$ for $\theta > \pi/4$.
The singularity is essential, not removable.}
\label{fig:path-dependence}
\end{figure}

\begin{itemize}[nosep]
\item Fixed $a > 0$, $b \to 0$: limit $= 1$.
\item $a = \lambda b$, $\lambda \in (0,1]$: limit $= \lambda$.
\item $a = b^\alpha$, $\alpha > 1$: limit $= 0$.
\end{itemize}
Every value in $[0,1]$ is achieved (Part~1; Lean:
\texttt{path\_dep\_fixed\_a},
\texttt{path\_dep\_proportional},
\texttt{path\_dep\_square},
\texttt{path\_any\_positive\_value}).
The essential singularity at the origin means the collapse, which
maps all of $(0,1)$ to $\tfrac{1}{2}$, cannot capture the
approach-direction information that determines the conditional.

\subsection{The Role of Dependence}

The min copula is only one joint; what about others?

Under the \emph{independence copula} ($\cred(A \land B) = a \cdot b$),
the conditional is $a \cdot b / b = a$ --- it does not depend on the
evidence at all.
The bridge trivially holds (collapsing $a$ gives the same answer
either way), but the conditional is uninformative.

For collapse($a$) = collapse($b$) = $\tfrac{1}{2}$ (both interior),
the set of min-copula conditionals is $(0, 1]$ --- maximally
underdetermined by the collapsed values
(Lean: \texttt{interior\_cond\_range}).
This is the precise sense in which the collapse loses too much
information: interior credences collapse to a single value, but
their ratios span almost the entire unit interval.

%=============================================================================
\section{Update Through the Bridge}
\label{sec:update}
%=============================================================================

The conditional bridge fails in general but holds at boundaries.
This section shows that Bayesian update inherits the bridge under
the same boundary conditions, defines qualitative update on
$\{0, \tfrac{1}{2}, 1\}$, and connects zero-evidence behavior
across the probabilistic and logical perspectives.

\subsection{Bayesian Update on $[0,1]$}

\begin{definition}[Bayesian Update; Part~1]
\label{def:bayesian}
Given a joint credence $\cred(A \land B)$ and evidence $\cred(B) > 0$,
the Bayesian posterior is
$\cred'(A) = \cred(A \land B) / \cred(B)$.
This satisfies the chain rule by construction:
$\cred'(A) \conjc \cred(B) = \cred(A \land B)$.
Lean: \texttt{bayesianUpdate\_chain\_rule},
\texttt{bayesianUpdate\_unique}.
\end{definition}

\begin{example}[Running Example: Bayesian Update]
\label{ex:bayesian}
Under the min copula:
$\cred(\text{rain} \land \text{wet}) = \min(0.8, 0.9) = 0.8$.
The Bayesian posterior:
$\cred'(\text{rain}) = 0.8/0.9 = 8/9 \approx 0.889$.
Rain was already likely; conditioning on wet streets (which are
correlated) makes it slightly more likely, but not certain.
\end{example}

\subsection{Qualitative Update on $\{0, \tfrac{1}{2}, 1\}$}

\begin{definition}[Qualitative Update; new]
\label{def:qual-update}
For $e, p \in \{0, \tfrac{1}{2}, 1\}$, the \emph{qualitative
update} is $\mathrm{qual}(e, p) = \mathrm{g\ddot{o}del}(e, p)$:
the G\"odel implication applied to evidence and prior.
Lean: \texttt{qualitativeUpdate}.
\end{definition}

\begin{center}
\small
\begin{tabular}{r|ccc}
\multicolumn{4}{c}{\textbf{Qualitative update table}:
$\mathrm{qual}(e, p) = \mathrm{g\ddot{o}del}(e, p)$} \\
$e \!\downarrow\;\; p$ & 0 & $\tfrac{1}{2}$ & 1 \\
\hline
0 & 1 & 1 & 1 \\
$\tfrac{1}{2}$ & 0 & 1 & 1 \\
1 & 0 & $\tfrac{1}{2}$ & 1
\end{tabular}
\normalsize
\end{center}
\smallskip
\noindent
Lean: \texttt{qualUpdate\_zero\_zero} through
\texttt{qualUpdate\_one\_one}.

The table has a natural reading.
Certain evidence ($e = 1$) preserves the prior.
Intermediate evidence ($e = \tfrac{1}{2}$) maps uncertain
conclusions to certainty (the G\"odel conditional's ``benefit of
the doubt'' for compatible evidence).
Zero evidence ($e = 0$) maps everything to $1$ --- the ex falso
row of the G\"odel conditional, which the chain rule avoids by
leaving the conditional unconstrained.

\subsection{Update Bridge Theorem}

\begin{theorem}[Update Bridge; new]
\label{thm:update-bridge}
Under boundary conditions (Theorem~\ref{thm:boundary}), collapsing
after min-copula conditioning equals qualitative update after
collapsing:
\[
\mathrm{collapse}\!\left(\frac{\min(a,b)}{b}\right)
= \mathrm{qual}(\mathrm{collapse}(b),\, \mathrm{collapse}(a)).
\]
That is: if $a = 0$, $a = 1$, $b = 1$, or $a \geq b$ (with
$b > 0$), then conditioning on $[0,1]$ followed by collapsing
gives the same result as collapsing followed by qualitative update.

\smallskip\noindent
Lean: \texttt{update\_bridge}.
\end{theorem}

\begin{proof}
Immediate from Theorem~\ref{thm:boundary}: the qualitative update
is defined as $\mathrm{g\ddot{o}del}(e, p)$, which is what the
boundary characterization gives.
\end{proof}

\begin{example}[Running Example: Update Bridge]
\label{ex:update-bridge}
Case (iv) from Theorem~\ref{thm:boundary}: suppose instead that
$\cred(\text{rain}) = 0.95 > 0.9 = \cred(\text{wet})$.
Then $\min(0.95, 0.9)/0.9 = 1$.
On the left:
$\mathrm{collapse}(1) = 1$.
On the right:
$\mathrm{qual}(\mathrm{collapse}(0.9), \mathrm{collapse}(0.95))
= \mathrm{qual}(\tfrac{1}{2}, \tfrac{1}{2})
= \mathrm{g\ddot{o}del}(\tfrac{1}{2}, \tfrac{1}{2}) = 1$.
Both paths agree.

But with the original values $\cred(\text{rain}) = 0.8 <
0.9 = \cred(\text{wet})$ (interior, $a < b$), the bridge fails:
$\mathrm{collapse}(8/9) = \tfrac{1}{2} \neq 1 =
\mathrm{qual}(\tfrac{1}{2}, \tfrac{1}{2})$.
\end{example}

\subsection{Jeffrey Update}

Jeffrey conditionalization handles uncertain evidence via a mixture:
$\cred'(A) = \condc{A}{B} \cdot q + \condc{A}{\negc B} \cdot (1-q)$,
where $q$ is the new credence for $B$~\cite{jeffrey1965}.
At $q = 1$ this reduces to Bayesian update; at $q = 0$ it conditions
on $\negc B$
(Lean: \texttt{jeffreyUpdate\_certain},
\texttt{jeffreyUpdate\_zero}).

For the Jeffrey update bridge to hold, \emph{both} conditional arms
must satisfy the boundary conditions of Theorem~\ref{thm:boundary}:
$\condc{A}{B}$ must be at a boundary, and so must
$\condc{A}{\negc B}$.
This is restrictive.
In the running example, unless $\cred(\text{rain})$ is $0$, $1$,
or at least as large as both $\cred(\text{wet})$ and
$\cred(\negc\text{wet})$, at least one arm is interior and the
Jeffrey bridge fails.

\subsection{Zero-Evidence Duality}
\label{sec:zero-duality}

The algebraic fact $c \cdot 0 = 0$ for all $c \in [0,1]$ is the
common root of three phenomena.

\begin{theorem}[Zero-Evidence Duality; new]
\label{thm:zero-duality}
The following three statements are simultaneously consequences of
the identity $c \conjc 0 = 0$:
\begin{enumerate}[label=\textup{(\roman*)},nosep]
\item \textbf{Underdetermined posteriors:}
  When evidence has credence $0$, any conditional
  $c \in [0,1]$ satisfies the chain rule.
  The set of consistent posteriors is $[0,1]$ --- the maximal
  imprecise posterior~\cite{walley1991,levi1980}.
\item \textbf{LP explosion failure:}
  The value $\tfrac{1}{2}$ makes both $A$ and $\negc A$
  LP-designated (since both $\tfrac{1}{2}$ and $\negc\tfrac{1}{2}
  = \tfrac{1}{2}$ are in $\{1, \tfrac{1}{2}\}$) without forcing
  any particular value for $B$.
  This is the witness that defeats explosion.
\item \textbf{Conditional bridge failure:}
  At evidence $0$, the chain rule is satisfied by every conditional
  value.
  No truth-functional $f$ can capture this:
  $f(\mathrm{collapse}(0), \mathrm{collapse}(a))$ must return a
  single value, but any value is equally valid.
\end{enumerate}

\smallskip\noindent
Lean: \texttt{zero\_evidence\_duality}
(packaging \texttt{zero\_evidence\_underdetermined},
\texttt{zero\_evidence\_paraconsistency},
\texttt{zero\_evidence\_bridge\_failure}).
\end{theorem}

\begin{proof}
(i)~$c \cdot 0 = 0$ for all $c$
(Lean: \texttt{conditioning\_zero\_any}).
(ii)~$v(A) = \tfrac{1}{2}$, $v(B) = 0$ satisfies
$v(A) > 0$ and $\negc v(A) = \tfrac{1}{2} > 0$ but $v(B) = 0$
(Lean: \texttt{lp\_no\_explosion}).
(iii)~For any $c_1, c_2 \in [0,1]$, both are consistent
conditionals at evidence $0$
(Lean: \texttt{conditioning\_zero\_any}).
\end{proof}

The duality provides a unified perspective on the paper's results.
The propositional bridge (Section~\ref{sec:bridge}) operates through
the connectives, which are homomorphic under collapse.
The conditional bridge (Section~\ref{sec:cond-bridge}) fails because
conditioning depends on \emph{ratios}, and the collapse loses ratio
information by mapping all of $(0,1)$ to $\tfrac{1}{2}$.
The zero-evidence case is the extreme: the ratio $0/0$ is maximally
underdetermined, which manifests simultaneously as paraconsistency
(the $\tfrac{1}{2}$ witness), imprecise probability (the vacuous
posterior), and bridge failure (no truth-functional conditional can
handle the underdetermination).

\begin{center}
\small
\begin{tabular}{lcccc}
\toprule
\textbf{System} & \textbf{Designated} & \textbf{Tautologies?}
  & \textbf{Excl.\ middle?} & \textbf{Explosion?} \\
\midrule
Classical & $\{1\}$ (two-valued) & Yes & Yes & Yes \\
K$_3$ & $\{1\}$ & No & No & Vacuously$^a$ \\
LP & $\{1, \tfrac{1}{2}\}$ & Yes & Yes & No \\
Cred: positivity & $\{c > 0\}$ & Yes & Yes & No$^b$ \\
Cred: certainty & $\{c = 1\}$ & No & No & Vacuously$^a$ \\
\bottomrule
\end{tabular}
\normalsize
\end{center}
\smallskip
\noindent
{\footnotesize
$^a$No valuation makes both $v(A) = 1$ and $\negc v(A) = 1$, so
explosion holds vacuously (premises jointly unsatisfiable).
$^b$The witness $v(A) = \tfrac{1}{2}$, $v(B) = 0$ refutes explosion
(Example~\ref{ex:explosion};
Lean: \texttt{lp\_no\_explosion}).}

%=============================================================================
\section{Conclusion}
\label{sec:conclusion}
%=============================================================================

Part~1 showed that the product De Morgan triplet on $[0,1]$
collapses to the Kleene lattice via a surjective homomorphism,
and that LP/K$_3$ consequence is equivalent to
positivity/certainty consequence at the formula level.
This paper asked whether the bridge extends to conditioning.

\textbf{The answer is no, in general.}
No truth-functional function on $\{0, \tfrac{1}{2}, 1\}$ can
serve as a conditional bridge
(Theorem~\ref{thm:no-cond-bridge}).
The impossibility is witnessed by two explicit pairs of interior
credences --- a constructive Lewis triviality for the product
De Morgan triplet.
The bridge holds exactly at boundaries: when the evidence is
certain, the conclusion is impossible or certain, or the conclusion
is at least as credible as the evidence
(Theorem~\ref{thm:boundary}).
Bayesian update inherits the bridge under these conditions
(Theorem~\ref{thm:update-bridge}).

\textbf{The common root is $c \cdot 0 = 0$.}
The algebraic identity that makes zero an annihilator for the
product simultaneously forces underdetermined posteriors at zero
evidence, provides the witness for LP explosion failure, and causes
the conditional bridge to break
(Theorem~\ref{thm:zero-duality}).
The propositional bridge operates on the connectives
$(\negc, \conjc, \disjc)$, which are preserved by the collapse.
The conditional bridge fails because conditioning depends on ratios,
and the collapse discards ratio information.

\textbf{Graded consequence.}
Graded consequence on $[0,1]$ with threshold $t$ exhibits a sharp
no-explosion threshold at $\tfrac{1}{2}$: for $t \leq \tfrac{1}{2}$,
the witness $v(A) = \tfrac{1}{2}$ defeats explosion; for
$t > \tfrac{1}{2}$, explosion holds vacuously (the premises are
unsatisfiable)
(Lean: \texttt{graded\_no\_explosion},
\texttt{graded\_explosion\_vacuous}).

\textbf{Open problems.}
\begin{enumerate}[nosep]
\item \textbf{Non-truth-functional conditional bridge.}
  A conditional bridge with an explicit dependence indicator
  (e.g., a function
  $g : \{0, \tfrac{1}{2}, 1\}^2 \times D \to
  \{0, \tfrac{1}{2}, 1\}$
  for some dependence type $D$) might recover the bridge at the
  cost of an additional parameter.
  What is the minimal such $D$?
\item \textbf{Four-valued FDE extension.}
  K$_3$ and LP are both fragments of Belnap--Dunn four-valued logic
  (FDE)~\cite{belnap1977,dunn1976}.
  Does the bridge theorem extend to FDE, and what is the
  corresponding consequence relation on $[0,1]$?
\item \textbf{Proof theory for chain-rule conditioning.}
  Is there a sound and complete calculus for chain-rule conditioning
  on $[0,1]$?
  H\'ajek's completeness theorem~\cite{hajek1998} answers this for
  product logic with residuation; replacing residuation with the
  chain rule changes the answer.
\item \textbf{Iterated update convergence.}
  Under what stochastic model for the evidence sequence does
  iterated Bayesian update drive $\cred_n(A) \to 1$, $0$, or
  $\tfrac{1}{2}$?
\end{enumerate}

\section*{Acknowledgements}

Omitted for blind review.

%=============================================================================
% Bibliography
%=============================================================================

\bibliographystyle{plain}
\bibliography{../cred}

%=============================================================================
% APPENDIX
%=============================================================================
\appendix

%=============================================================================
\section{Preliminaries}
\label{app:prelim}
%=============================================================================

This appendix collects standard definitions; see Part~1 for full
treatment.

\textbf{T-norms.}
A \emph{t-norm} is a binary operation
$T : [0,1]^2 \to [0,1]$ that is commutative, associative,
monotone, and satisfies $T(a, 1) = a$ for all
$a$~\cite{klement2000}.
The three basic continuous t-norms are product ($a \cdot b$),
G\"odel ($\min(a,b)$), and \L ukasiewicz ($\max(a+b-1, 0)$).

\textbf{Copulas.}
A copula $C : [0,1]^2 \to [0,1]$ satisfies: grounded
($C(a,0) = C(0,b) = 0$), unit marginals ($C(a,1) = a$, $C(1,b) = b$),
and 2-increasing~\cite{nelsen2006}.
The Fr\'echet-Hoeffding bounds apply:
$\max(a+b-1, 0) \leq C(a,b) \leq \min(a,b)$.

\textbf{Residuation.}
Given a left-continuous t-norm $T$, the residuum is
$a \to_T b = \sup\{c : T(a,c) \leq b\}$~\cite{hajek1998,galatos2007}.
For the product: $a \to_P b = \min(b/a, 1)$ when $a > 0$, and
$0 \to_P b = 1$ (ex falso).

\textbf{Coherence.}
De Finetti~\cite{definetti1937} introduced coherence: an assignment
is coherent if no finite combination of bets yields a sure loss.
The chain rule
$\cred(A \mid B) \cdot \cred(B) = \cred(A \land B)$
is necessary for coherence~\cite{coletti2002,gilio2014}.

%=============================================================================
\section{Lean Formalization}
\label{app:lean}
%=============================================================================

The results are formalized across seven Lean~4 modules (with
Mathlib), totaling 3581 lines with zero \texttt{sorry}:

\begin{itemize}[nosep]
\item \texttt{Cred/Basic.lean}: core algebra, conditioning, fixed
  points, collapse, Bayes consistency, path dependence, congruence
  building blocks (1804 lines).
\item \texttt{Cred/Congruence.lean}: two-level congruence
  classification --- \texttt{UnitCongruence}, \texttt{RealCongruence},
  Kleene witness (396 lines).
\item \texttt{Cred/Consequence.lean}: K$_3$/LP designation,
  formula evaluation, collapse-evaluation commutativity,
  bridge theorems, structural rules,
  graded consequence, no-explosion (505 lines).
\item \texttt{Cred/Valuation.lean}: valuations, collapse composition,
  Fr\'echet bounds (175 lines).
\item \texttt{Cred/Update.lean}: Bayesian/Jeffrey update, chain-rule
  preservation (138 lines).
\item \texttt{Cred/Predicate.lean}: graded predicates,
  quantifiers, Russell fixed point (210 lines).
\item \texttt{Cred/CondBridge.lean}: conditional bridge impossibility,
  boundary characterization, qualitative update, update bridge,
  zero-evidence duality (353 lines).
\end{itemize}

\textbf{Key definitions (this paper):}
\begin{center}
\small
\begin{tabular}{ll}
\toprule
\textbf{Concept} & \textbf{Lean name} \\
\midrule
Min-copula conditional & \texttt{Cred.minCopulaCond} \\
Qualitative update & \texttt{Cred.qualitativeUpdate} \\
G\"odel implication & \texttt{Cred.ThreeVal.godel\_impl} \\
Product-residuated implication & \texttt{Cred.ThreeVal.prod\_resid\_impl} \\
Propositional formula & \texttt{Cred.Formula} \\
Credence evaluation & \texttt{Cred.evalCred} \\
Three-valued evaluation & \texttt{Cred.evalThreeVal} \\
Bayesian update & \texttt{Cred.bayesianUpdate} \\
Jeffrey update & \texttt{Cred.jeffreyUpdate} \\
\bottomrule
\end{tabular}
\normalsize
\end{center}

\textbf{Key theorems (this paper):}
\begin{center}
\small
\begin{tabular}{ll}
\toprule
\textbf{Paper result} & \textbf{Lean name} \\
\midrule
Collapse-eval commutativity (Prop.~\ref{prop:collapse-eval})
  & \texttt{collapse\_eval\_eq}, \texttt{lift\_eval\_eq} \\
LP formula bridge (Thm.~\ref{thm:lp-formula-bridge})
  & \texttt{lp\_formula\_bridge} \\
K$_3$ formula bridge (Thm.~\ref{thm:k3-formula-bridge})
  & \texttt{k3\_formula\_bridge} \\
No cond.\ bridge (Thm.~\ref{thm:no-cond-bridge})
  & \texttt{no\_truthfunctional\_cond\_bridge} \\
Boundary characterization (Thm.~\ref{thm:boundary})
  & \texttt{cond\_bridge\_boundary} \\
Interior failure (Thm.~\ref{thm:interior-fail})
  & \texttt{cond\_bridge\_fails\_interior} \\
Interior conditional range
  & \texttt{interior\_cond\_range} \\
Qualitative update table
  & \texttt{qualUpdate\_*} \\
Update bridge (Thm.~\ref{thm:update-bridge})
  & \texttt{update\_bridge} \\
Zero-evidence duality (Thm.~\ref{thm:zero-duality})
  & \texttt{zero\_evidence\_duality} \\
Bayesian update (Def.~\ref{def:bayesian})
  & \texttt{bayesianUpdate\_chain\_rule/unique} \\
Jeffrey update
  & \texttt{jeffreyUpdate\_certain/zero} \\
LP no explosion (Ex.~\ref{ex:explosion})
  & \texttt{lp\_no\_explosion} \\
K$_3$ no tautology
  & \texttt{k3\_no\_tautology} \\
Graded no-explosion
  & \texttt{graded\_no\_explosion} \\
Graded explosion vacuous
  & \texttt{graded\_explosion\_vacuous} \\
G\"odel $=$ prod.\ resid.
  & \texttt{godel\_impl\_eq\_prod\_resid} \\
\bottomrule
\end{tabular}
\normalsize
\end{center}

Verification: \texttt{cd lean \&\& lake build}.
Toolchain: Lean 4.16.0 with Mathlib 4.16.0.

\end{document}
